\documentclass[a4paper]{report}
\usepackage{hyperref}
\usepackage{booktabs}
\begin{document}

\chapter{Table row by row with bold \& italic example}

\noindent
\begin{tabular}{lll}
\toprule
\textbf{City} & \textbf{Country} & \textbf{Size} \\
\midrule
Mariehamn & Finland & Small \\
\textit{Copenhagen} & \textit{Denmark} & \textit{Large} \\
\textit{\textbf{Tokyo}} & \textit{\textbf{Japan}} & \textit{\textbf{Huge}} \\
\bottomrule
\end{tabular}
\par\medskip

\noindent
\begin{tabular}{lll}
\toprule
\textit{Capital} & \textit{Country} & \textit{Continent} \\
\midrule
Oslo & Norway & Europe \\
Tokyo & Japan & Asia \\
\textbf{Berlin} & \textbf{Germany} & \textbf{Europe} \\
Kairo & Egypt & Africa \\
\textit{\textbf{Brussels}} & \textit{\textbf{Belgium}} & \textit{\textbf{Europe}} \\
\bottomrule
\end{tabular}
\par\medskip

Note: As the rows are added and written one by one, the library cannot know the width of columns in future rows, this will make markdown formatted output look a bit odd when reading the markdown file, however any further formatting from markdown will hide this and make the table look as expected.

\end{document}
